\documentclass[UTF8]{ctexart}
\usepackage{xeCJK}
\usepackage{geometry}
\geometry{left=2cm,right=2cm,top=2.5cm,bottom=2.5cm}
\setlength{\headheight}{12.7pt}

\usepackage{titlesec}
\titleformat{\section}
  {\normalfont\large\bfseries}
  {\thesection}
  {1em}
  {}
\titlespacing*{\section}
  {0pt}
  {3.5ex plus 1ex minus .2ex}
  {2.3ex plus .2ex}

\usepackage{amsmath}
\usepackage{amssymb}
\usepackage{amsthm}
\usepackage{enumitem}
\setlist[enumerate]{itemsep=0pt,topsep=0pt,parsep=0pt,leftmargin=0pt,itemindent=2.5em}
\setlist[itemize]{itemsep=0pt,topsep=0pt,parsep=0pt,leftmargin=0pt,itemindent=2.5em}
\usepackage{fancyhdr}
\pagestyle{fancy}
\fancyhf{}
\usepackage{graphicx}
\usepackage{hyperref}
\usepackage{indentfirst}
\usepackage{lmodern}


\begin{document}
\setlength{\abovedisplayskip}{1pt}
\setlength{\belowdisplayskip}{1pt}

\section*{类}
\addcontentsline{toc}{section}{类}

直观上可以把满足某性质的集合的全体看作一个对象, 但是有时这样的对象在集合论语言内不存在,
在元语言中称之为``类''.

\vspace{10pt}

\hypertarget{sec:类}{}
\noindent\textbf{类}

任取公式$\varphi(p_1,\cdots,p_n,x)$,
它可以指定一个以$p_1,\cdots,p_n$为参数的类$\mathbf{A}$, 在集合论语言中
\vspace{3pt}
\[
    x\in\mathbf{A}\overset{\mathrm{def}}{\leftrightarrow}\varphi.
\]

特别地, 集合也可以视为一个类.

\vspace{10pt}

下面省略对参数的说明, 默认有$p_1,\cdots,p_n$ (也可以没有)作为参数.

\vspace{10pt}

\hypertarget{sec:真类}{}
\noindent\textbf{真类}

给定类$\mathbf{A}$. 如果
$\neg\,\exists a\,\forall x:\,x\in\mathbf{A}\leftrightarrow x\in a$,
就称$\mathbf{A}$是一个真类.

\vspace{10pt}

\hypertarget{sec:子类}{}
\noindent\textbf{子类}

给定类$\mathbf{A},\mathbf{B}$. 如果$\forall x:x\in\mathbf{A}\to x\in\mathbf{B}$,
就称$\mathbf{A}$是$\mathbf{B}$的子类, 记$\mathbf{A}\subset\mathbf{B}$.

进一步, 如果$\forall x:x\in\mathbf{A}\leftrightarrow x\in\mathbf{B}$,
就记$\mathbf{A}=\mathbf{B}$.

\vspace{10pt}

\hypertarget{sec:类的二元运算}{}
\noindent\textbf{类的二元运算}

给定类$\mathbf{A},\mathbf{B}$,

定义它们的二元并$\mathbf{A}\cup\mathbf{B}$:
$x\in\mathbf{A}\cup\mathbf{B}\overset{\mathrm{def}}{\leftrightarrow}
x\in\mathbf{A}\lor x\in\mathbf{B}$;

定义它们的二元交$\mathbf{A}\cap\mathbf{B}$:
$x\in\mathbf{A}\cap\mathbf{B}\overset{\mathrm{def}}{\leftrightarrow}
x\in\mathbf{A}\land x\in\mathbf{B}$;

定义它们的二元积$\mathbf{A}\times\mathbf{B}$:
$p\in\mathbf{A}\times\mathbf{B}\overset{\mathrm{def}}{\leftrightarrow}
\exists x\in\mathbf{A}\,\exists y\in\mathbf{B}:p=(x,y)$.

\vspace{10pt}

\hypertarget{sec:类的一般运算}{}
\noindent\textbf{类的一般运算}

给定类$\mathbf{A}$.

定义它的一般并$\displaystyle\bigcup\mathbf{A}$:
$\displaystyle x\in\bigcup\mathbf{A}\overset{\mathrm{def}}{\leftrightarrow}
\exists y\in\mathbf{A}:x\in y$.

定义它的一般交$\displaystyle\bigcap\mathbf{A}$:
$\displaystyle x\in\bigcap\mathbf{A}\overset{\mathrm{def}}{\leftrightarrow}
\forall y\in\mathbf{A}:x\in y$.

\vspace{10pt}

\hypertarget{sec:集合宇宙}{}
\noindent\textbf{集合宇宙}

公式$x=x$指定一个没有参数的类, 记作$\mathbf{V}$, 称作集合宇宙.
显然$x\in\mathbf{V}$始终成立.





\newpage

\section*{类映射}
\addcontentsline{toc}{section}{类映射}

类映射是映射的推广, 只关注对应关系, 不关注整体对象.

\vspace{10pt}

\hypertarget{sec:类映射}{}
\noindent\textbf{类映射}

给定类$\mathbf{A},\mathbf{B}$和$\mathbf{P}\subset\mathbf{A}\times\mathbf{B}$,
如果
\[
    \forall x\in\mathbf{A}\,\exists!y\in\mathbf{B}:(x,y)\in\mathbf{P},
\]
那么称$\mathbf{P}$是从$\mathbf{A}$到$\mathbf{B}$的类映射. 把$y$记为$\mathbf{P}(x)$.

\vspace{10pt}

\hypertarget{sec:真类到集合没有单射}{}
\noindent\textbf{真类到集合没有单射}

给定类$\mathbf{A}$, 集合$b$和$\mathbf{A}$上的类映射$\mathbf{P}$, 则以下矛盾:
\begin{enumerate}
    \item $\mathbf{A}$是真类,
    \item 对任意的$x\in\mathbf{A}$, $\mathbf{P}(x)\in b$,
    \item 对任意的$x,x'\in\mathbf{A}$,
    如果$\mathbf{P}(x)=\mathbf{P}(x')$, 那么$x=x'$.
\end{enumerate}

\noindent{\kaishu 证明.}

使用\href{01 ZFC 公理#nameddest=sec:分离公理模式}{分离公理}, 存在$c$使得
\[
    \forall y:\quad y\in c\leftrightarrow y\in b\land\exists x:y=\mathbf{P}(x),
\]
从而$\forall y\in c\,\exists! x:y=\mathbf{P}(x)$.

再使用\href{01 ZFC 公理#nameddest=sec:替换公理模式}{替换公理}, 存在$a$使得
\[
    \forall x:\quad x\in a\leftrightarrow\exists y\in c:y=\mathbf{P}(x),
\]
同时显然
\[
    \forall x:\quad x\in\mathbf{A}\leftrightarrow\exists y\in c:y=\mathbf{P}(x),
\]
从而$\forall x:x\in\mathbf{A}\leftrightarrow x\in a$, 与$\mathbf{A}$是真类矛盾.

\end{document}